% ============================================================
%  GUÍA 1: Fundamentos de Álgebra
%  Curso de Introducción a la Universidad — Precálculo y Cálculo I
%  Autor: Ing. Gabriel Astudillo
%  Clase cubierta: 1
% ============================================================

\documentclass[12pt, a4paper]{article}

% ── Paquetes ─────────────────────────────────────────────────
\usepackage{mathwizards-guia}
\mwguia{Guía 1 — Fundamentos de Álgebra}{Ing. Gabriel Astudillo $\cdot$ Curso Preuniversitario}

\begin{document}

% ── Portada / Título ─────────────────────────────────────────
\mwportada{Guía de Estudio 1}{Fundamentos de Álgebra}{Exponentes, Racionalización, Factorización, Fracciones Algebraicas, Ecuaciones y Desigualdades}

\vspace{4pt}
\begin{cuadroinfo}
\small
\textbf{Presentación:} Esta primera guía sienta las bases operativas y conceptuales de todo el curso preuniversitario y de cálculo diferencial. Dominar las leyes de los exponentes, la racionalización de diversos órdenes, todos los casos de factorización, el álgebra de fracciones complejas, la resolución de ecuaciones y el trabajo con desigualdades es indispensable para abordar con éxito los temas de funciones, límites y derivadas. Cada sección incluye un marco teórico exhaustivo, ejercicios resueltos paso a paso y ejercicios propuestos con su clave de respuestas.
\\[4pt]
\textbf{Nivel de Dificultad:}\quad
\facil{} Básico / Directo \quad
\medio{} Intermedio \quad
\dificil{} Desafiante \quad
\muyDificil{} Muy Difícil / Reto
\end{cuadroinfo}

\vspace{8pt}

\section{Leyes de Exponentes, Radicales y Racionalización}
% ============================================================

\subsection{Leyes fundamentales de exponentes y radicales}

\begin{cuadroteoria}
\textbf{Leyes de los exponentes:} Para $a, b > 0$ y $m, n \in \mathbb{R}$:
\begin{center}
\small
\begin{tabular}{lll}
\toprule
\textbf{Ley} & \textbf{Fórmula} & \textbf{Ejemplo} \\
\midrule
Producto de potencias de igual base & $a^m \cdot a^n = a^{m+n}$ & $x^3 \cdot x^5 = x^8$ \\
Cociente de potencias de igual base & $\dfrac{a^m}{a^n} = a^{m-n} \quad (a \neq 0)$ & $\dfrac{x^7}{x^2} = x^5$ \\
Potencia de una potencia & $(a^m)^n = a^{m \cdot n}$ & $(x^3)^4 = x^{12}$ \\
Potencia de un producto & $(ab)^n = a^n \cdot b^n$ & $(2x)^3 = 8x^3$ \\
Potencia de un cociente & $\left(\dfrac{a}{b}\right)^n = \dfrac{a^n}{b^n} \quad (b \neq 0)$ & $\left(\dfrac{x}{3}\right)^2 = \dfrac{x^2}{9}$ \\
Exponente cero & $a^0 = 1 \quad (a \neq 0)$ & $(-5)^0 = 1$ \\
Exponente negativo & $a^{-n} = \dfrac{1}{a^n} \quad (a \neq 0)$ & $x^{-3} = \dfrac{1}{x^3}$ \\
Fracción con exponente negativo & $\left(\dfrac{a}{b}\right)^{-n} = \left(\dfrac{b}{a}\right)^n$ & $\left(\dfrac{2}{3}\right)^{-2} = \dfrac{9}{4}$ \\
Exponente fraccionario & $a^{m/n} = \sqrt[n]{a^m} = (\sqrt[n]{a})^m$ & $8^{2/3} = (\sqrt[3]{8})^2 = 4$ \\
\bottomrule
\end{tabular}
\end{center}
\end{cuadroteoria}

\begin{cuadroadvertencia}
\textbf{Errores comunes a evitar:}
\begin{itemize}[leftmargin=*]
  \item $a^m \cdot a^n \neq a^{m \cdot n}$ \quad (los exponentes se \emph{suman}, no se multiplican).
  \item $(a + b)^n \neq a^n + b^n$ \quad (la potencia \textbf{no} es distributiva respecto a la suma ni a la resta).
  \item $a^{-n}$ \textbf{no} es un número negativo: representa el recíproco positivo $\frac{1}{a^n}$.
  \item $\sqrt{a^2} = |a|$ para todo $a \in \mathbb{R}$. Por ejemplo, $\sqrt{(-4)^2} = |-4| = 4 \neq -4$.
\end{itemize}
\end{cuadroadvertencia}

\newpage

\subsection{Técnicas de Racionalización}

La racionalización consiste en transformar una expresión con radicales en el numerador o denominador en una fracción equivalente libre de radicales en dicha posición.

\begin{cuadropasos}
\textbf{1. Denominador Monomio:}
\begin{itemize}[leftmargin=*]
  \item \textbf{Raíz cuadrada $\sqrt{a}$:} Se multiplica numerador y denominador por $\sqrt{a}$:
        $$\frac{k}{\sqrt{a}} \cdot \frac{\sqrt{a}}{\sqrt{a}} = \frac{k\sqrt{a}}{a}$$
  \item \textbf{Raíz de orden $n$ ($\sqrt[n]{a^k}$ con $k < n$):} Se multiplica por $\sqrt[n]{a^{n-k}}$ para completar el exponente:
        $$\frac{k}{\sqrt[n]{a^m}} \cdot \frac{\sqrt[n]{a^{n-m}}}{\sqrt[n]{a^{n-m}}} = \frac{k\sqrt[n]{a^{n-m}}}{a}$$
        Ejemplo para raíz cúbica: $\dfrac{k}{\sqrt[3]{a}} \cdot \dfrac{\sqrt[3]{a^2}}{\sqrt[3]{a^2}} = \dfrac{k\sqrt[3]{a^2}}{a}$.
\end{itemize}

\textbf{2. Denominador Binomio con Raíces Cuadradas (Conjugado de Orden 2):}
Se utiliza la diferencia de cuadrados $(A + B)(A - B) = A^2 - B^2$:
$$\frac{k}{\sqrt{a} \pm \sqrt{b}} \cdot \frac{\sqrt{a} \mp \sqrt{b}}{\sqrt{a} \mp \sqrt{b}} = \frac{k(\sqrt{a} \mp \sqrt{b})}{a - b}$$

\textbf{3. Denominador Binomio con Raíces Cúbicas (Conjugado de Orden 3):}
Se utilizan las identidades de suma y diferencia de cubos:
\begin{align*}
  (\sqrt[3]{a} + \sqrt[3]{b})(\sqrt[3]{a^2} - \sqrt[3]{ab} + \sqrt[3]{b^2}) &= a + b \\
  (\sqrt[3]{a} - \sqrt[3]{b})(\sqrt[3]{a^2} + \sqrt[3]{ab} + \sqrt[3]{b^2}) &= a - b
\end{align*}
Multiplicando por el factor trinomio adecuado, se eliminan las raíces cúbicas del denominador.

\textbf{4. Racionalización del Numerador (Cocientes Incrementales de Límites):}
Multiplicar por el conjugado del numerador para eliminar la indeterminación $\frac{0}{0}$.
\end{cuadropasos}

\subsection{Ejercicios resueltos: Exponentes y Racionalización}

\textbf{Ejemplo R1.} Simplifica dejando solo exponentes positivos: $\dfrac{(3x^2 y^{-3})^3 \cdot (x^{-2} y^4)^2}{9x^{-1} y^2}$.

\begin{cuadrosolucion}
\textbf{Solución:}
\begin{enumerate}
  \item Aplicamos potencias de potencias:
        $$\frac{27 x^6 y^{-9} \cdot x^{-4} y^8}{9 x^{-1} y^2} = \frac{27 x^{6 - 4} y^{-9 + 8}}{9 x^{-1} y^2} = \frac{27 x^2 y^{-1}}{9 x^{-1} y^2}$$
  \item Dividimos coeficientes y restamos exponentes:
        $$\frac{27}{9} \cdot x^{2 - (-1)} \cdot y^{-1 - 2} = 3 x^3 y^{-3} = \frac{3x^3}{y^3}.$$
\end{enumerate}
\end{cuadrosolucion}

\newpage

\textbf{Ejemplo R2.} Racionaliza el denominador monomio con raíz cúbica: $\dfrac{5}{\sqrt[3]{2x}}$.

\begin{cuadrosolucion}
\textbf{Solución:}
Multiplicamos por $\sqrt[3]{(2x)^2} = \sqrt[3]{4x^2}$ tanto arriba como abajo:
$$\frac{5}{\sqrt[3]{2x}} \cdot \frac{\sqrt[3]{4x^2}}{\sqrt[3]{4x^2}} = \frac{5\sqrt[3]{4x^2}}{\sqrt[3]{(2x)^3}} = \frac{5\sqrt[3]{4x^2}}{2x}.$$
\end{cuadrosolucion}

\textbf{Ejemplo R3.} Racionaliza el denominador binomio: $\dfrac{6}{\sqrt{7} - 2}$.

\begin{cuadrosolucion}
\textbf{Solución:}
Multiplicamos por el conjugado del denominador $\sqrt{7} + 2$:
$$\frac{6}{\sqrt{7} - 2} \cdot \frac{\sqrt{7} + 2}{\sqrt{7} + 2} = \frac{6(\sqrt{7} + 2)}{(\sqrt{7})^2 - 2^2} = \frac{6(\sqrt{7} + 2)}{7 - 4} = \frac{6(\sqrt{7} + 2)}{3} = 2(\sqrt{7} + 2) = 2\sqrt{7} + 4.$$
\end{cuadrosolucion}

\textbf{Ejemplo R4.} Racionaliza el denominador que contiene una raíz cúbica en binomio: $\dfrac{4}{\sqrt[3]{x} + 2}$.

\begin{cuadrosolucion}
\textbf{Solución:}
Identificamos $A = \sqrt[3]{x}$ y $B = 2$. Usamos $A^3 + B^3 = (A + B)(A^2 - AB + B^2)$ con factor racionalizante $A^2 - AB + B^2 = \sqrt[3]{x^2} - 2\sqrt[3]{x} + 4$:
$$\frac{4}{\sqrt[3]{x} + 2} \cdot \frac{\sqrt[3]{x^2} - 2\sqrt[3]{x} + 4}{\sqrt[3]{x^2} - 2\sqrt[3]{x} + 4} = \frac{4(\sqrt[3]{x^2} - 2\sqrt[3]{x} + 4)}{(\sqrt[3]{x})^3 + 2^3} = \frac{4(\sqrt[3]{x^2} - 2\sqrt[3]{x} + 4)}{x + 8}.$$
\end{cuadrosolucion}

\textbf{Ejemplo R5.} Racionaliza el numerador de $\dfrac{\sqrt{x + 9} - 3}{x}$ y simplifica para $x \neq 0$.

\begin{cuadrosolucion}
\textbf{Solución:}
Multiplicamos por el conjugado del numerador:
$$\frac{(\sqrt{x + 9} - 3)(\sqrt{x + 9} + 3)}{x(\sqrt{x + 9} + 3)} = \frac{(x + 9) - 9}{x(\sqrt{x + 9} + 3)} = \frac{x}{x(\sqrt{x + 9} + 3)} = \frac{1}{\sqrt{x + 9} + 3}.$$
\end{cuadrosolucion}

\subsection{Ejercicios propuestos: Exponentes y Racionalización}

\begin{enumerate}[leftmargin=*]
  \item \facil{} Simplifica usando leyes de potencias:
        $$a)\; (2x^3 y^{-2})^4 \qquad\qquad b)\; \frac{(x^2 y^3)^2}{(x y^2)^3} \qquad\qquad c)\; \left(\frac{3a^{-2} b}{a^3 b^{-1}}\right)^{-2}$$

  \item \facil{} Escribe como una única potencia fraccionaria irreducible:
        $$a)\; \sqrt{x\sqrt{x}} \qquad\qquad b)\; \sqrt[3]{x^2 \sqrt[4]{x}} \qquad\qquad c)\; \frac{\sqrt[3]{a^2}}{\sqrt{a}}$$

  \item \facil{} Racionaliza el denominador monomio:
        $$a)\; \frac{4}{\sqrt{6}} \qquad\qquad b)\; \frac{10}{\sqrt{5x}} \qquad\qquad c)\; \frac{3}{\sqrt[3]{4}}$$

  \item \medio{} Racionaliza el denominador binomio con raíces cuadradas:
        $$a)\; \frac{5}{\sqrt{3} + \sqrt{2}} \qquad\qquad b)\; \frac{x - 4}{\sqrt{x} - 2} \qquad\qquad c)\; \frac{12}{3 - \sqrt{5}}$$

  \item \medio{} Racionaliza el denominador monomio de orden superior:
        $$a)\; \frac{6}{\sqrt[4]{8}} \qquad\qquad b)\; \frac{x}{\sqrt[5]{x^3}} \qquad\qquad c)\; \frac{2}{\sqrt[3]{9x^2}}$$

  \item \medio{} Racionaliza el numerador:
        $$a)\; \frac{\sqrt{x + 4} - 2}{x} \qquad\qquad b)\; \frac{\sqrt{x + h} - \sqrt{x}}{h} \qquad\qquad c)\; \frac{\sqrt{2x + 1} - 3}{x - 4}$$

  \item \dificil{} Racionaliza el denominador binomio con raíces cúbicas:
        $$a)\; \frac{2}{\sqrt[3]{5} - \sqrt[3]{3}} \qquad\qquad b)\; \frac{1}{\sqrt[3]{x} - 1} \qquad\qquad c)\; \frac{6}{\sqrt[3]{9} - \sqrt[3]{3} + 1}$$

  \item \dificil{} Racionaliza el numerador con raíz cúbica:
        $$\frac{\sqrt[3]{x + 8} - 2}{x}$$
\end{enumerate}

\newpage

% ============================================================

\section{Factorización Sistemática Caso por Caso}
% ============================================================

\subsection{Catálogo completo de casos de factorización}

La factorización es el proceso inverso al desarrollo de productos notables; transforma sumas y restas algebraicas en productos de factores irreducibles.

\begin{cuadroteoria}
\textbf{Resumen de Casos de Factorización:}
\begin{enumerate}[leftmargin=*]
  \item \textbf{Factor Común Monomio y Polinomio:}
        $$ab + ac = a(b + c) \qquad\qquad x(a + b) + y(a + b) = (a + b)(x + y)$$
  \item \textbf{Factorización por Agrupación de Términos (4 o 6 términos):}
        $$ax + ay + bx + by = a(x + y) + b(x + y) = (x + y)(a + b)$$
  \item \textbf{Diferencia de Cuadrados:}
        $$a^2 - b^2 = (a - b)(a + b)$$
  \item \textbf{Trinomio Cuadrado Perfecto (TCP):}
        $$a^2 \pm 2ab + b^2 = (a \pm b)^2$$
  \item \textbf{Trinomio de la forma $x^2 + bx + c$:}
        $$x^2 + (p + q)x + pq = (x + p)(x + q) \quad \text{donde } p+q = b \text{ y } p \cdot q = c$$
  \item \textbf{Trinomio de la forma $ax^2 + bx + c$ ($a \neq 1$):}
        Se descompone $bx$ en dos números que multipliquen $a \cdot c$ y sumen $b$, para luego aplicar agrupación.
  \item \textbf{Suma y Diferencia de Cubos:}
        $$a^3 + b^3 = (a + b)(a^2 - ab + b^2) \qquad\qquad a^3 - b^3 = (a - b)(a^2 + ab + b^2)$$
  \item \textbf{Completación de Cuadrados (Método de Sophie Germain):}
        Sumar y restar un término intermedio para formar un TCP y luego una diferencia de cuadrados.
\end{enumerate}
\end{cuadroteoria}

\begin{cuadroadvertencia}
\textbf{Estrategia de ataque en factorización combinada:}
1. Extraer siempre primero el máximo factor común (MFC). \\
2. Contar el número de términos restantes:
   \begin{itemize}
     \item 2 términos: ¿Es diferencia de cuadrados ($a^2 - b^2$) o suma/diferencia de cubos ($a^3 \pm b^3$)?
     \item 3 términos: ¿Es TCP, $x^2 + bx + c$, o $ax^2 + bx + c$?
     \item 4 o más términos: ¿Agrupación $2+2$, agrupación $3+1$ (TCP menos cuadrado)?
   \end{itemize}
3. Factorizar completamente hasta que todos los factores sean irreducibles en $\mathbb{R}$.
\end{cuadroadvertencia}

\newpage

\subsection{Ejercicios resueltos: Casos de Factorización}

\textbf{Ejemplo R6 (Factor Común y Agrupación).} Factoriza: $a)\; 6x^3 y^2 - 9x^2 y^4 + 15x^4 y^3 \qquad b)\; x^3 - 3x^2 - 4x + 12$.

\begin{cuadrosolucion}
\textbf{Solución:}
\begin{enumerate}
  \item[a)] El máximo factor común es $3x^2 y^2$:
        $$6x^3 y^2 - 9x^2 y^4 + 15x^4 y^3 = 3x^2 y^2(2x - 3y^2 + 5x^2 y).$$
  \item[b)] Agrupamos de dos en dos:
        $$(x^3 - 3x^2) - (4x - 12) = x^2(x - 3) - 4(x - 3) = (x - 3)(x^2 - 4) = (x - 3)(x - 2)(x + 2).$$
\end{enumerate}
\end{cuadrosolucion}

\textbf{Ejemplo R7 (Trinomio $ax^2 + bx + c$).} Factoriza: $6x^2 - 11x - 10$.

\begin{cuadrosolucion}
\textbf{Solución:}
Multiplicamos $a \cdot c = 6 \cdot (-10) = -60$. Buscamos dos números que multipliquen $-60$ y sumen $-11$: son $-15$ y $+4$.
Reescribimos el término central y agrupamos:
$$6x^2 - 15x + 4x - 10 = 3x(2x - 5) + 2(2x - 5) = (2x - 5)(3x + 2).$$
\end{cuadrosolucion}

\textbf{Ejemplo R8 (Suma y Diferencia de Cubos).} Factoriza: $a)\; 27x^3 + 64 \qquad b)\; 8x^6 - y^3$.

\begin{cuadrosolucion}
\textbf{Solución:}
\begin{enumerate}
  \item[a)] $27x^3 + 64 = (3x)^3 + 4^3 = (3x + 4)((3x)^2 - (3x)(4) + 4^2) = (3x + 4)(9x^2 - 12x + 16)$.
  \item[b)] $8x^6 - y^3 = (2x^2)^3 - y^3 = (2x^2 - y)(4x^4 + 2x^2 y + y^2)$.
\end{enumerate}
\end{cuadrosolucion}

\textbf{Ejemplo R9 (Completación de Cuadrados / Sophie Germain).} Factoriza: $x^4 + 4y^4$.

\begin{cuadrosolucion}
\textbf{Solución:}
Para formar un trinomio cuadrado perfecto necesitamos el término intermedio $2(x^2)(2y^2) = 4x^2 y^2$. Sumamos y restamos $4x^2 y^2$:
$$x^4 + 4x^2 y^2 + 4y^4 - 4x^2 y^2 = (x^2 + 2y^2)^2 - (2xy)^2$$
Aplicamos diferencia de cuadrados:
$$= (x^2 + 2y^2 - 2xy)(x^2 + 2y^2 + 2xy) = (x^2 - 2xy + 2y^2)(x^2 + 2xy + 2y^2).$$
\end{cuadrosolucion}

\textbf{Ejemplo R10 (Combinación Múltiple).} Factoriza completamente: $x^6 - 64$.

\begin{cuadrosolucion}
\textbf{Solución:}
Se debe aplicar primero diferencia de cuadrados y luego suma/diferencia de cubos:
$$x^6 - 64 = (x^3)^2 - 8^2 = (x^3 - 8)(x^3 + 8) = (x - 2)(x^2 + 2x + 4)(x + 2)(x^2 - 2x + 4).$$
\end{cuadrosolucion}

\newpage

\subsection{Ejercicios propuestos: Factorización}

\begin{enumerate}[leftmargin=*, resume]
  \item \facil{} \textbf{Factor Común Monomio y Polinomio:}
        $$a)\; 12x^4 y^3 - 18x^3 y^5 \qquad b)\; 4a^2 b - 8ab^2 + 12a^3 b^3 \qquad c)\; 2x(a - 1) - y(a - 1)$$

  \item \facil{} \textbf{Diferencia de Cuadrados:}
        $$a)\; 49x^2 - 36 \qquad\qquad b)\; 100a^4 - 81b^6 \qquad\qquad c)\; (x + 2)^2 - 25$$

  \item \facil{} \textbf{Trinomio Cuadrado Perfecto:}
        $$a)\; x^2 + 14x + 49 \qquad\qquad b)\; 9x^2 - 30xy + 25y^2 \qquad\qquad c)\; 4a^2 + 12ab + 9b^2$$

  \item \medio{} \textbf{Factorización por Agrupación de Términos:}
        $$a)\; 2x^3 - 3x^2 + 4x - 6 \qquad b)\; ax + bx - ay - by \qquad c)\; 3x^3 + 2x^2 - 27x - 18$$

  \item \medio{} \textbf{Trinomios $x^2 + bx + c$ y $ax^2 + bx + c$:}
        $$a)\; x^2 - 8x + 12 \qquad b)\; x^2 + 2x - 35 \qquad c)\; 3x^2 + 7x + 2 \qquad d)\; 4x^2 - 4x - 15$$

  \item \medio{} \textbf{Suma y Diferencia de Cubos:}
        $$a)\; x^3 + 125 \qquad\qquad b)\; 64x^3 - 1 \qquad\qquad c)\; 8a^3 + 27b^6$$

  \item \dificil{} \textbf{Factorización de Casos Especiales (Completación de Cuadrados):}
        $$a)\; a^4 + 64b^4 \qquad\qquad b)\; x^4 + x^2 + 1 \qquad\qquad c)\; x^4 - 7x^2 y^2 + y^4$$

  \item \dificil{} \textbf{Factorización Combinada Avanzada:}
        $$a)\; 2x^4 - 32 \qquad b)\; x^3 - 2x^2 - x + 2 \qquad c)\; x^4 - 13x^2 + 36 \qquad d)\; 4x^2 - y^2 + 4x + 1$$
\end{enumerate}

\newpage

% ============================================================

\section{Fracciones Algebraicas y Expresiones Complejas}
% ============================================================

\subsection{Operaciones con Fracciones Algebraicas}

Una fracción algebraica es un cociente de polinomios $\frac{P(x)}{Q(x)}$ con $Q(x) \neq 0$.

\begin{cuadroteoria}
\textbf{Operaciones Fundamentales:}
\begin{itemize}[leftmargin=*]
  \item \textbf{Simplificación:} Se factorizan completamente numerador y denominador y se cancelan los factores comunes:
        $$\frac{P(x) \cdot K(x)}{Q(x) \cdot K(x)} = \frac{P(x)}{Q(x)} \quad (K(x) \neq 0)$$
  \item \textbf{Multiplicación:} $\dfrac{A}{B} \cdot \dfrac{C}{D} = \dfrac{A \cdot C}{B \cdot D}$ \quad (factorizar antes de multiplicar para simplificar).
  \item \textbf{División:} $\dfrac{A}{B} \div \dfrac{C}{D} = \dfrac{A}{B} \cdot \dfrac{D}{C} = \dfrac{A \cdot D}{B \cdot C}$.
  \item \textbf{Suma y Resta (Homogéneas):} $\dfrac{A}{M} \pm \dfrac{B}{M} = \dfrac{A \pm B}{M}$.
  \item \textbf{Suma y Resta (Heterogéneas):} Se factorizan los denominadores para hallar el \textbf{mínimo común múltiplo (m.c.m.)} de los denominadores:
        $$\frac{A}{B} + \frac{C}{D} = \frac{A \cdot \frac{\text{mcm}}{B} + C \cdot \frac{\text{mcm}}{D}}{\text{mcm}}$$
\end{itemize}
\end{cuadroteoria}

\begin{cuadropasos}
\textbf{Simplificación de Fracciones Complejas (Fracciones dentro de Fracciones):}
\begin{itemize}[leftmargin=*]
  \item \textbf{Método 1 (m.c.m. global):} Multiplicar el numerador principal y el denominador principal por el m.c.m. de todos los denominadores secundarios.
  \item \textbf{Método 2 (Suma separada):} Reducir el numerador a una sola fracción, el denominador a una sola fracción, y aplicar la división de extremos y medios (doble C):
        $$\frac{\frac{A}{B}}{\frac{C}{D}} = \frac{A \cdot D}{B \cdot C}$$
\end{itemize}
\end{cuadropasos}

\newpage

\subsection{Ejercicios resueltos: Fracciones Algebraicas}

\textbf{Ejemplo R11 (Simplificación).} Simplifica completamente: $\dfrac{x^3 - 8}{x^2 - 4}$.

\begin{cuadrosolucion}
\textbf{Solución:}
Factorizamos numerador (diferencia de cubos) y denominador (diferencia de cuadrados):
$$\frac{x^3 - 8}{x^2 - 4} = \frac{(x - 2)(x^2 + 2x + 4)}{(x - 2)(x + 2)} = \frac{x^2 + 2x + 4}{x + 2} \quad (x \neq 2).$$
\end{cuadrosolucion}

\textbf{Ejemplo R12 (Suma y Resta con m.c.m.).} Efectúa y simplifica: $\dfrac{2}{x^2 - 1} - \dfrac{1}{x^2 - x}$.

\begin{cuadrosolucion}
\textbf{Solución:}
\begin{enumerate}
  \item Factorizamos denominadores: $x^2 - 1 = (x - 1)(x + 1)$ y $x^2 - x = x(x - 1)$.
  \item El m.c.m. es $x(x - 1)(x + 1)$.
  \item Convertimos a denominador común:
        $$\frac{2x - 1(x + 1)}{x(x - 1)(x + 1)} = \frac{2x - x - 1}{x(x - 1)(x + 1)} = \frac{x - 1}{x(x - 1)(x + 1)} = \frac{1}{x(x + 1)}.$$
\end{enumerate}
\end{cuadrosolucion}

\textbf{Ejemplo R13 (Multiplicación y División).} Efectúa: $\dfrac{x^2 - 9}{x^2 + 4x + 4} \div \dfrac{2x - 6}{x + 2}$.

\begin{cuadrosolucion}
\textbf{Solución:}
Invertimos la segunda fracción y factorizamos todos los términos:
$$\frac{(x - 3)(x + 3)}{(x + 2)^2} \cdot \frac{x + 2}{2(x - 3)} = \frac{(x + 3)(x - 3)(x + 2)}{2(x + 2)^2 (x - 3)} = \frac{x + 3}{2(x + 2)}.$$
\end{cuadrosolucion}

\textbf{Ejemplo R14 (Fracción Compleja).} Simplifica: $\dfrac{\frac{1}{x + h} - \frac{1}{x}}{h}$.

\begin{cuadrosolucion}
\textbf{Solución:}
Operamos el numerador con m.c.m. $x(x + h)$:
$$\frac{\frac{x - (x + h)}{x(x + h)}}{h} = \frac{\frac{-h}{x(x + h)}}{h} = \frac{-h}{h \cdot x(x + h)} = -\frac{1}{x(x + h)} \quad (h \neq 0).$$
\end{cuadrosolucion}

\newpage

\subsection{Ejercicios propuestos: Fracciones Algebraicas}

\begin{enumerate}[leftmargin=*, resume]
  \item \facil{} \textbf{Simplificación directa:}
        $$a)\; \frac{x^2 - 9}{x^2 - 5x + 6} \qquad\qquad b)\; \frac{2x^2 + 6x}{x^2 - 9} \qquad\qquad c)\; \frac{x^2 + 6x + 9}{x^3 + 27}$$

  \item \medio{} \textbf{Multiplicación y División:}
        $$a)\; \frac{x^2 - 4}{x^2 + 3x} \cdot \frac{x^2 - 9}{x^2 + 4x + 4} \qquad\qquad b)\; \frac{2x^2 - 5x - 3}{x^2 - 1} \div \frac{4x^2 - 1}{x + 1}$$

  \item \medio{} \textbf{Suma y Resta de Fracciones:}
        $$a)\; \frac{3}{x - 2} - \frac{2}{x + 3} \qquad\qquad b)\; \frac{x}{x^2 - 4} + \frac{1}{x + 2} \qquad\qquad c)\; \frac{2x}{x^2 - 1} - \frac{1}{x - 1} + \frac{1}{x + 1}$$

  \item \dificil{} \textbf{Fracciones Complejas (Fracción de Fracciones):}
        $$a)\; \frac{\frac{1}{x} + \frac{1}{y}}{\frac{1}{x^2} - \frac{1}{y^2}} \qquad\qquad b)\; \frac{1 - \frac{2}{x} - \frac{3}{x^2}}{1 - \frac{9}{x^2}} \qquad\qquad c)\; \frac{\frac{x + 1}{x - 1} - \frac{x - 1}{x + 1}}{\frac{1}{x - 1} + \frac{1}{x + 1}}$$

  \item \dificil{} \textbf{Operación Combinada Integral:}
        $$\left(\frac{x + 1}{x - 1} - \frac{x - 1}{x + 1}\right) \cdot \frac{x^2 - 1}{4x} + \frac{1}{x + 2}$$
\end{enumerate}

\newpage

% ============================================================

\section{Ecuaciones Lineales, Cuadráticas y con Valor Absoluto}
% ============================================================

\subsection{Resolución de ecuaciones}

\begin{cuadroteoria}
\textbf{Ecuación lineal:} $ax + b = 0$ con $a \neq 0$. Solución: $x = -\dfrac{b}{a}$.

\textbf{Ecuación cuadrática:} $ax^2 + bx + c = 0$ con $a \neq 0$. Se resuelve por:
\begin{itemize}[leftmargin=*]
  \item \emph{Factorización}: si $ax^2 + bx + c = (px + q)(rx + s) = 0 \implies x = -\frac{q}{p}$ o $x = -\frac{s}{r}$.
  \item \emph{Fórmula general}: $x = \dfrac{-b \pm \sqrt{b^2 - 4ac}}{2a}$.
  \item \emph{Completando cuadrados}.
\end{itemize}
El discriminante $D = b^2 - 4ac$ determina la naturaleza de las soluciones:
\begin{itemize}
  \item $D > 0$: dos soluciones reales y distintas.
  \item $D = 0$: una solución real (raíz doble).
  \item $D < 0$: sin soluciones reales (dos complejas conjugadas).
\end{itemize}

\textbf{Ecuación con valor absoluto:} $|X| = k$ con $k \geq 0$ equivale a $X = k$ o $X = -k$.
\end{cuadroteoria}

\subsection{Ejercicios resueltos: Ecuaciones}

\textbf{Ejemplo R15.} Resuelve $\dfrac{3x - 1}{2} = \dfrac{x + 5}{4}$.

\begin{cuadrosolucion}
\textbf{Solución:}
\begin{enumerate}
  \item Multiplicamos ambos lados por 4: $2(3x - 1) = x + 5 \implies 6x - 2 = x + 5$.
  \item Agrupamos términos semejantes: $5x = 7 \implies x = \dfrac{7}{5}$.
\end{enumerate}
\end{cuadrosolucion}

\textbf{Ejemplo R16.} Resuelve $2x^2 + 3x - 2 = 0$.

\begin{cuadrosolucion}
\textbf{Solución:}
Por fórmula general con $a = 2, b = 3, c = -2$:
$$x = \frac{-3 \pm \sqrt{9 - 4(2)(-2)}}{4} = \frac{-3 \pm \sqrt{25}}{4} = \frac{-3 \pm 5}{4} \implies x = \frac{1}{2} \quad\text{o}\quad x = -2.$$
\end{cuadrosolucion}

\textbf{Ejemplo R17.} Resuelve $|2x - 5| = 7$.

\begin{cuadrosolucion}
\textbf{Solución:}
Equivale a dos casos lineales:
\begin{itemize}
  \item $2x - 5 = 7 \implies 2x = 12 \implies x = 6$.
  \item $2x - 5 = -7 \implies 2x = -2 \implies x = -1$.
\end{itemize}
Soluciones: $x = 6$ o $x = -1$.
\end{cuadrosolucion}

\newpage

\subsection{Ejercicios propuestos: Ecuaciones}

\begin{enumerate}[leftmargin=*, resume]
  \item \facil{} Resuelve: $a)\; 5x - 8 = 3x + 10 \qquad\qquad b)\; \dfrac{2x}{3} - \dfrac{x}{4} = 5$

  \item \facil{} Resuelve: $a)\; x^2 - 49 = 0 \qquad\qquad b)\; x^2 + 6x + 9 = 0 \qquad\qquad c)\; x^2 - 3x - 10 = 0$

  \item \medio{} Resuelve por fórmula general: $3x^2 - 5x + 1 = 0$

  \item \medio{} Resuelve ecuaciones con valor absoluto:
        $$a)\; |3x + 2| = 11 \qquad\qquad b)\; |x - 4| = |2x + 1|$$

  \item \medio{} Resuelve la ecuación racional: $\dfrac{x + 2}{x - 1} = \dfrac{2x}{x + 3}$

  \item \dificil{} Resuelve la ecuación bicuadrada: $x^4 - 13x^2 + 36 = 0$

  \item \dificil{} Resuelve con verificación de raíces extrañas: $\sqrt{x + 3} = x - 3$
\end{enumerate}

\newpage

% ============================================================

\section{Desigualdades e Intervalos}
% ============================================================

\subsection{Propiedades y notación de intervalos}

\begin{cuadroteoria}
\textbf{Propiedades de las desigualdades:} Para $a, b, c \in \mathbb{R}$:
\begin{itemize}[leftmargin=*]
  \item Si $a < b$ y $b < c$, entonces $a < c$ (transitividad).
  \item Si $a < b$, entonces $a + c < b + c$ (sumar una constante conserva el orden).
  \item Si $a < b$ y $c > 0$, entonces $ac < bc$ (multiplicar por positivo conserva el orden).
  \item Si $a < b$ y \textbf{$c < 0$}, entonces \textbf{$ac > bc$} (multiplicar por negativo \textbf{invierte} el sentido).
\end{itemize}

\textbf{Desigualdades con Valor Absoluto ($k > 0$):}
\begin{itemize}[leftmargin=*]
  \item $|X| \le k \iff -k \le X \le k \iff X \in [-k, k]$
  \item $|X| \ge k \iff X \le -k \quad\text{o}\quad X \ge k \iff X \in (-\infty, -k] \cup [k, \infty)$
\end{itemize}
\end{cuadroteoria}

\subsection{Ejercicios resueltos: Desigualdades}

\textbf{Ejemplo R18.} Resuelve $x^2 - x - 6 > 0$.

\begin{cuadrosolucion}
\textbf{Solución:}
Factorizamos $(x - 3)(x + 2) > 0$. Los puntos críticos son $x = -2$ y $x = 3$.
\begin{center}
\small
\begin{tabular}{lccc}
\toprule
\textbf{Intervalo} & $(-\infty, -2)$ & $(-2, 3)$ & $(3, \infty)$ \\
\midrule
Signo de $(x - 3)(x + 2)$ & $+$ & $-$ & $+$ \\
\bottomrule
\end{tabular}
\end{center}
Solución: $x \in (-\infty, -2) \cup (3, \infty)$.
\end{cuadrosolucion}

\textbf{Ejemplo R19.} Resuelve $\dfrac{x - 1}{x + 2} \geq 0$.

\begin{cuadrosolucion}
\textbf{Solución:}
Restricción: $x \neq -2$. Raíz del numerador: $x = 1$.
Puntos críticos en la recta: $-2$ y $1$.
El cociente es positivo o cero en: $(-\infty, -2) \cup [1, \infty)$. (El $-2$ va abierto por división entre cero).
\end{cuadrosolucion}

\newpage

\subsection{Ejercicios propuestos: Desigualdades}

\begin{enumerate}[leftmargin=*, resume]
  \item \facil{} Resuelve y expresa en intervalos: $a)\; 3x + 4 > x + 10 \qquad\qquad b)\; -2x + 5 \le 11$

  \item \facil{} Resuelve: $a)\; x^2 - 9 < 0 \qquad\qquad b)\; x^2 - 4x \ge 0$

  \item \medio{} Resuelve: $2x^2 + 3x - 5 < 0$

  \item \medio{} Resuelve con valor absoluto: $a)\; |x + 1| \le 4 \qquad\qquad b)\; |2x - 3| > 5$

  \item \medio{} Resuelve la desigualdad fraccionaria: $\dfrac{2x - 1}{x + 3} \le 1$

  \item \dificil{} Resuelve analizando por casos: $|x - 2| + |x + 3| < 9$

  \item \dificil{} Resuelve la desigualdad racional cuadrática: $\dfrac{x^2 - 5x + 4}{x^2 - 1} \ge 0$
\end{enumerate}

\newpage

% ============================================================

% ──────────────────────────────────────────────────────────
%  NUEVOS EJERCICIOS PROPUESTOS (15) — INSERCIÓN MANUAL
% ──────────────────────────────────────────────────────────
\subsection{Ejercicios propuestos: Repaso Integral de Álgebra I}

\begin{enumerate}[leftmargin=*, resume]
  \item \facil{} Simplifica dejando solo exponentes positivos:
        $$(3x^2 y^{-3})^{-2} \cdot (2x^{-1} y^2)^3$$

  \item \facil{} Racionaliza los denominadores:
        $$a)\; \frac{8}{\sqrt{12}} \qquad\qquad\qquad b)\; \frac{5}{\sqrt[3]{16x}}$$

  \item \medio{} Factoriza completamente en $\mathbb{R}$:
        $$x^4 - 5x^2 + 4$$

  \item \medio{} Simplifica la fracción compleja (indica las restricciones de $x$ e $y$):
        $$\frac{\dfrac{1}{x} - \dfrac{1}{y}}{\dfrac{1}{x^2} - \dfrac{1}{y^2}}$$

  \item \medio{} Resuelve la ecuación con valor absoluto:
        $$|2x - 3| = |x + 4|$$

  \item \dificil{} Resuelve la ecuación racional y verifica que no existan raíces extrañas:
        $$\frac{x}{x + 1} - \frac{2}{x - 1} = 1$$

  \item \dificil{} Resuelve la ecuación con radicales:
        $$\sqrt{x + 3} + \sqrt{x - 2} = 5$$

  \item \muyDificil{} Resuelve la desigualdad racional analizando por puntos críticos:
        $$\frac{x^2 - 4x + 3}{x^2 + 2x - 3} \ge 0$$
\end{enumerate}

\subsection{Ejercicios propuestos: Repaso Integral de Álgebra II}

\begin{enumerate}[leftmargin=*, resume]
  \item \facil{} Simplifica dejando solo exponentes positivos:
        $$\frac{(a^3 b^{-2})^2 (2a b^3)^3}{4 a^2 b^{-1}}$$

  \item \medio{} Racionaliza el numerador y simplifica para $x \neq 4$:
        $$\frac{\sqrt{x + 5} - 3}{x - 4}$$

  \item \medio{} Factoriza completamente el polinomio y resuelve la ecuación $P(x) = 0$:
        $$P(x) = x^3 + 2x^2 - 9x - 18$$

  \item \dificil{} Resuelve la ecuación mediante cambio de variable:
        $$(x^2 + x + 1)^2 - 7(x^2 + x + 1) + 10 = 0$$

  \item \dificil{} Resuelve la desigualdad con valor absoluto analizando por casos:
        $$|x - 1| + |x + 2| \le 5$$

  \item \muyDificil{} Simplifica la fracción compleja anidada (indica las restricciones):
        $$\frac{1 + \dfrac{1}{1 + \dfrac{1}{x}}}{1 - \dfrac{1}{1 + \dfrac{1}{x}}}$$

  \item \medio{} \textbf{Aplicación:} El largo de un rectángulo excede en 4 m a su ancho y su área es de $96\ \text{m}^2$. Halla sus dimensiones.
\end{enumerate}

% ============================================================
\section{Clave de Respuestas}
% ============================================================

\subsection*{Sección 1: Exponentes y Racionalización (Ejercicios 1--8)}
\begin{enumerate}[leftmargin=*, label=\textbf{\arabic*.}]
  \item a) $\dfrac{16x^{12}}{y^8}$ \quad b) $x y^3$ \quad c) $\dfrac{a^{10}}{9 b^4}$
  \item a) $x^{3/4}$ \quad b) $x^{11/12}$ \quad c) $a^{1/6}$
  \item a) $\dfrac{2\sqrt{6}}{3}$ \quad b) $\dfrac{2\sqrt{5x}}{x}$ \quad c) $\dfrac{3\sqrt[3]{2}}{2}$
  \item a) $5(\sqrt{3} - \sqrt{2})$ \quad b) $\sqrt{x} + 2$ \quad c) $3(3 + \sqrt{5})$
  \item a) $3\sqrt[4]{2}$ \quad b) $\sqrt[5]{x^2}$ \quad c) $\dfrac{2\sqrt[3]{3x}}{3x}$
  \item a) $\dfrac{1}{\sqrt{x + 4} + 2}$ \quad b) $\dfrac{1}{\sqrt{x + h} + \sqrt{x}}$ \quad c) $\dfrac{2}{\sqrt{2x + 1} + 3}$
  \item a) $\sqrt[3]{25} + \sqrt[3]{15} + \sqrt[3]{9}$ \quad b) $\dfrac{\sqrt[3]{x^2} + \sqrt[3]{x} + 1}{x - 1}$ \quad c) $2(\sqrt[3]{3} + 1)$
  \item $\dfrac{1}{\sqrt[3]{(x + 8)^2} + 2\sqrt[3]{x + 8} + 4}$
\end{enumerate}


\subsection*{Sección 2: Factorización (Ejercicios 9--16)}
\begin{enumerate}[leftmargin=*, label=\textbf{\arabic*.}]
  \setcounter{enumi}{8}
  \item a) $6x^3 y^3(2x - 3y^2)$ \quad b) $4ab(a - 2b + 3a^2 b^2)$ \quad c) $(a - 1)(2x - y)$
  \item a) $(7x - 6)(7x + 6)$ \quad b) $(10a^2 - 9b^3)(10a^2 + 9b^3)$ \quad c) $(x - 3)(x + 7)$
  \item a) $(x + 7)^2$ \quad b) $(3x - 5y)^2$ \quad c) $(2a + 3b)^2$
  \item a) $(2x - 3)(x^2 + 2)$ \quad b) $(x - y)(a + b)$ \quad c) $(3x + 2)(x - 3)(x + 3)$
  \item a) $(x - 2)(x - 6)$ \quad b) $(x + 7)(x - 5)$ \quad c) $(3x + 1)(x + 2)$ \quad d) $(2x + 3)(2x - 5)$
  \item a) $(x + 5)(x^2 - 5x + 25)$ \quad b) $(4x - 1)(16x^2 + 4x + 1)$ \quad c) $(2a + 3b^2)(4a^2 - 6ab^2 + 9b^4)$
  \item a) $(a^2 - 4ab + 8b^2)(a^2 + 4ab + 8b^2)$ \quad b) $(x^2 - x + 1)(x^2 + x + 1)$ \quad c) $(x^2 - 3xy - y^2)(x^2 + 3xy - y^2)$
  \item a) $2(x - 2)(x + 2)(x^2 + 4)$ \quad b) $(x - 1)(x + 1)(x - 2)$ \quad c) $(x - 2)(x + 2)(x - 3)(x + 3)$ \quad d) $(2x + 1 - y)(2x + 1 + y)$
\end{enumerate}


\subsection*{Sección 3: Fracciones Algebraicas (Ejercicios 17--21)}
\begin{enumerate}[leftmargin=*, label=\textbf{\arabic*.}]
  \setcounter{enumi}{16}
  \item a) $\dfrac{x + 3}{x - 2}$ \quad b) $\dfrac{2x}{x - 3}$ \quad c) $\dfrac{x + 3}{x^2 - 3x + 9}$
  \item a) $\dfrac{(x - 2)(x - 3)}{x(x + 2)}$ \quad b) $\dfrac{x - 3}{2x - 1}$
  \item a) $\dfrac{x + 13}{(x - 2)(x + 3)}$ \quad b) $\dfrac{2(x - 1)}{(x - 2)(x + 2)}$ \quad c) $\dfrac{2}{x^2 - 1}$
  \item a) $\dfrac{xy}{y - x}$ \quad b) $\dfrac{x + 1}{x + 3}$ \quad c) $\dfrac{2x}{x^2 + 1}$
  \item $1 + \dfrac{1}{x + 2} = \dfrac{x + 3}{x + 2}$
\end{enumerate}


\subsection*{Sección 4: Ecuaciones (Ejercicios 22--28)}
\begin{enumerate}[leftmargin=*, label=\textbf{\arabic*.}]
  \setcounter{enumi}{21}
  \item a) $x = 9$ \quad b) $x = 12$
  \item a) $x = \pm 7$ \quad b) $x = -3$ (doble) \quad c) $x = 5$ o $x = -2$
  \item $x = \dfrac{5 \pm \sqrt{13}}{6}$
  \item a) $x = 3$ o $x = -\dfrac{13}{3}$ \quad b) $x = -5$ o $x = 1$
  \item $x = \dfrac{7 \pm \sqrt{73}}{2}$
  \item $x = \pm 2$ o $x = \pm 3$
  \item $x = 6$ ($x = 1$ es extraña y se descarta)
\end{enumerate}


\subsection*{Sección 5: Desigualdades (Ejercicios 29--35)}
\begin{enumerate}[leftmargin=*, label=\textbf{\arabic*.}]
  \setcounter{enumi}{28}
  \item a) $(3, \infty)$ \quad b) $[-3, \infty)$
  \item a) $(-3, 3)$ \quad b) $(-\infty, 0] \cup [4, \infty)$
  \item $\left(-\dfrac{5}{2}, 1\right)$
  \item a) $[-5, 3]$ \quad b) $(-\infty, -1) \cup (4, \infty)$
  \item $(-3, 4]$
  \item $(-5, 4)$
  \item $(-\infty, -1) \cup [4, \infty)$
\end{enumerate}


% ──────────────────────────────────────────────────────────
%  NUEVAS RESPUESTAS (15) — INSERCIÓN MANUAL
% ──────────────────────────────────────────────────────────
\subsection*{Sección 6: Repaso Integral de Álgebra I (Ejercicios 36--43)}
\begin{enumerate}[leftmargin=*, label=\textbf{\arabic*.}]
  \setcounter{enumi}{35}
  \item $\dfrac{8y^{12}}{9x^{7}}$
  \item a) $\dfrac{4\sqrt{3}}{3}$ \quad b) $\dfrac{5\sqrt[3]{4x^2}}{4x}$
  \item $(x - 1)(x + 1)(x - 2)(x + 2)$
  \item $\dfrac{xy}{x + y}$ (con $x \neq \pm y$, $x, y \neq 0$)
  \item $x = 7$ \quad o \quad $x = -\dfrac{1}{3}$
  \item $x = -\dfrac{1}{3}$ (verifica: los denominadores no se anulan)
  \item $x = 6$ (comprobación: $\sqrt{9} + \sqrt{4} = 5$)
  \item $(-\infty, -3) \cup [3, \infty)$ \quad (con $x \neq 1$; el punto $x = 1$ da la forma $0/0$ y se excluye)
\end{enumerate}

\subsection*{Sección 7: Repaso Integral de Álgebra II (Ejercicios 44--50)}
\begin{enumerate}[leftmargin=*, label=\textbf{\arabic*.}]
  \setcounter{enumi}{43}
  \item $2a^7 b^6$
  \item $\dfrac{1}{\sqrt{x + 5} + 3}$
  \item $(x + 2)(x - 3)(x + 3)$; \quad raíces: $x = -2$, $x = 3$, $x = -3$
  \item Con $u = x^2 + x + 1$: $u = 2$ o $u = 5$. Soluciones:
        $x = \dfrac{-1 \pm \sqrt{5}}{2}$ \quad y \quad $x = \dfrac{-1 \pm \sqrt{17}}{2}$ \quad (4 soluciones reales)
  \item $[-3, 2]$
  \item $2x + 1$ \quad (con $x \neq 0$ y $x \neq -1$)
  \item Ancho $8$ m, largo $12$ m \quad (de $x^2 + 4x - 96 = 0$; se descarta $x = -12$)
\end{enumerate}
\end{document}
